% Template for post or presentation abstracts submitted to ILAS 2022
% 1. Fill in the presenter's information (name, affiliation, email
% address) and talk information (title of presentation, abstract).
% 2. Compile to make sure there are no errors.
% 3. Save the .tex file for your abstract with a distinctive name,
% ideally "FamilynameGivenname.tex".
% 4. Upload your .tex file to https://clr.ie/129557
%       

\documentclass[a4paper]{amsart} 
\usepackage{amssymb} 
\usepackage{hyperref}
\pagestyle{empty}
\date{} 

%% IGNORE THE NEXT 4 LINES
\makeatletter % 
\def\labelenumi{\theenumi.}
\def\theenumi{\@arabic\c@enumi}
\makeatother
%% STOP IGNORING NOW

\begin{document}

\title{Perfect 1-factorisations and Hamiltonian Latin squares}
\author{Ian Wanless} 
\address{Monash University} 
\email{ian.wanless@monash.edu}

\maketitle

  A 1-\emph{factorisation} of a graph is a decomposition of the edges of the graph into 1-factors (perfect matchings). The 1-factorisation is \emph{perfect} if the union of any two of its 1-factors is a Hamilton cycle.  A P1F of the complete bipartite graph $K_{n,n}$ is equivalent to a \emph{row-Hamiltonian Latin square} of order $n$. These are Latin squares with no non-trivial Latin subrectangles; equivalently, the permutation which maps any row to any other row is an $n$-cycle. Each Latin square has six \emph{conjugates} obtained by uniformly permuting its (row, column, symbol) triples. Let $\nu(L)$ denote the number of conjugates of $L$ that are row-Hamiltonian. It is easy to see that $\nu(L)\in\{0,2,4,6\}$ and that $\nu=0$ can be achieved for all $n>3$. At the other extreme, $\nu=6$ is achieved by the so-called \emph{atomic} Latin squares, including the Cayley tables of cyclic groups of prime order. There is also a known infinite family with $\nu=2$. By converting our problem into linear algebra, we were able to find the first infinite family with $\nu=4$. We can build Latin squares in which \emph{every} pair of rows form a Hamilton cycle and \emph{no} pair of columns form a Hamilton cycle. As a corollary, we answer a question on quasigroup varieties posed by Falconer in 1970.

%% Coauthor details here (optional - delete if not applicable)
% and funding acknowledgment (optional - delete if not applicable)
\emph{This is joint work with Jack Allsop (Monash)}

\end{document}


% Please leave the next bit alone  
%%% Local Variables: 
%%% mode: latex
%%% TeX-master: "../ILAS-2022"
%%% End:
